% MJ81--2C
% Akutes Lungenödem, schwer
% 14.0
% 2013-11-30

\label{m:lungenoedem-schwer}
\EE{Taktik}: \Speech{\Ca{Vitale Bedrohung}
Belastung minimieren, rasche
medikamentöse Therpie (veranlassen); ärztliche Therapie
notwendig}

\begin{itemize}
  \item \TxMassVitMK
  \begin{itemize}
    \item Lagerung: Oberkörper hoch, \E{wenn möglich Beine
    herabhängen lassen}

    Striktes \EE{Bewegungsverbot!}
    \item  Notarzt!
    \item Reanimationsbereitschaft!
  \end{itemize}
  \item Beengende Kleidung öffnen
  \item \Z{Evtl. Gabe von Nitroglycerin-Spray durch einen
  Notfallsanitäter}
  \item Transportentscheidung: Internistische Intensiv- bzw. Überwachungstation
\end{itemize}

\Cave{Patienten mit Lungenödem und brodelndem Atemgeräusch oder ABCD-Problem 
sind grundsätzlich notarztpflichtig. 

Auch bei kurzer Transportzeit hat die Stabilisierung des Patienten vor Ort Vorrang.

Bereits der Transport in das Fahrzeug kann gefährlich sein!}

\Customized{
\begin{description}
  \item[\CustAsboe] Beim kardialen Lungenödem ist die Gabe
  von Nitrolingual Pumpspray gem. Algorithmus durch NFS vorgesehen.
  \cite{Asb921Memo-AL-2011-000001}.
\end{description}